\Askhsh{13687}{2}
\begin{ylh}{1}
    \item 3.2. 1ο Κριτήριο ισότητας τριγώνων
    \item 4.2. Τέμνουσα δύο ευθειών - Ευκλείδειο αίτημα
    \item 4.6. Άθροισμα γωνιών τριγώνου.
\end{ylh}
% ===================================================================
%  Αρχείο: 13687.tex (Fragment)
%  Αυτόματη μετατροπή μέσω doc2latex.py
% ===================================================================

ΘΕΜΑ 2

Σε κύκλο με κέντρο Ο και ακτίνα R θεωρούμε επίκεντρη γωνία Α$\widehat{\text{Ο}}$Β$\  = \ $40\textsuperscript{ο}. Προεκτείνουμε τις ακτίνες ΟΑ και ΟΒ κατά τμήματα ΑΓ και ΒΔ αντίστοιχα, έτσι ώστε ΑΓ$\  = \ $ΟΑ και ΒΔ$\  = \ $ΟΒ, όπως φαίνεται στο παρακάτω σχήμα.

\begin{alist}
\item Να αποδείξετε ότι Ο$\widehat{\text{Α}}$Β$\  = \ $Ο$\widehat{\text{Β}}$Α$\  = 7\text{0}\text{ο}$.

\monades{10}

\item Να υπολογίσετε τις γωνίες $\text{O}\widehat{\text{Γ}}$Δ και $\text{O}\widehat{\text{Δ}}$Γ.

\monades{10}

\item Να αποδείξετε ότι ΑΒ//ΓΔ.

\monades{5}

\includesvg[width=2.7141in,height=3.74016in]{/home/spyros/Μαθηματικά/Trapeza_Thematwn/A_Lykeiou/Geometry/extracted/13687/media/image2.svg}
\end{alist}
